% swift-attributes-directives.tex — a catalogue of Swift attributes, modifiers and
% compiler directives: library evolution (@inlinable,
% @usableFromInline, @frozen, @backDeployed, @export, @inline), API surface
% (@discardableResult, @available forms, @main, @testable, @_spi, package, @retroactive),
% ObjC interop (@objc, @objcMembers, @nonobjc, dynamic), concurrency markers
% (@preconcurrency, @unchecked Sendable, @TaskLocal, closure attributes, nonisolated),
% ownership (consuming/borrowing/consume/copy), #if conditions, #warning/#error,
% #sourceLocation, the magic literals (#file/#fileID/#filePath/#function/#dsohandle),
% #isolation, #externalMacro — with a diagram of what crosses a library-evolution boundary.
% Deliberately NOT repeated (see Related): access levels and @inlinable basics
% (swift-keypaths-access-control), @frozen enums + @unknown default
% (swift-enums-pattern-matching), macro roles (macros-and-result-builders), isolation
% semantics (swift6-strict-concurrency), ~Copyable (ownership-and-memory-layout).
% @_spi, @inline(__always) are underscored/unofficial and marked so.
% Sources (checked 2026-10-06) — Swift Evolution, every number and version against
% download.swift.org/swift-evolution/v1/evolution.json: SE-0193
% (@inlinable/@usableFromInline), SE-0260 (@frozen), SE-0376 (@backDeployed, 5.8),
% SE-0496 (@inline(always), 6.3: errors when a direct call can't be inlined; implies
% @inlinable for public), SE-0497 (@export(implementation|interface), 6.3), SE-0281
% (@main), SE-0383 (@UIApplicationMain deprecated, error in Swift 6, 5.10), SE-0386
% (package, 5.9), SE-0364 (@retroactive, 6.0), SE-0337 (@preconcurrency, 5.6), SE-0423
% (@preconcurrency conformances + dynamic isolation checks, 6.0), SE-0311 (task locals),
% SE-0377 (borrowing/consuming + copy, 5.9), SE-0366 (consume, 5.9), SE-0274 (#fileID /
% #filePath; #file = #fileID in Swift 6 mode), SE-0362 (hasFeature, 5.8), SE-0367
% (hasAttribute, 5.8), SE-0308 (#if around postfix members, 5.5), SE-0196 (#warning /
% #error), SE-0420 (#isolation, 6.0), SE-0382 (#externalMacro); The Swift Programming
% Language — Attributes, Statements ("Compiler Control Statements": inactive branches are
% parsed except under swift()/compiler(); targetEnvironment = simulator | macCatalyst).
% @labels: area=swift kind=concept platform=apple topic=language,build discussion=151
% @tags: attributes, library-evolution, inlinable, usablefrominline, frozen, backdeployed, available, objcmembers, preconcurrency, tasklocal, conditional-compilation, magic-literals
\documentclass[8pt]{extarticle}
\usepackage{printup-sheet}
\usepackage{array}
\usepackage{tabularx}

\lstdefinelanguage{SwiftSheet}{
  morekeywords={protocol,class,final,struct,enum,func,var,let,init,case,
    if,else,return,self,nil,true,false,in,public,internal,static,import,
    macro,some,any,isolated},
  sensitive=true, morecomment=[l]{//}, morestring=[b]",
  literate={->}{{\hbox{-}\hbox{>}}}2 {==}{{\hbox{=}\hbox{=}}}2 {!=}{{\hbox{!}\hbox{=}}}2
           {??}{{\hbox{?}\hbox{?}}}2 {...}{{\hbox{.}\hbox{.}\hbox{.}}}3
           {..<}{{\hbox{.}\hbox{.}\hbox{<}}}3 {>=}{{\hbox{>}\hbox{=}}}2
           {<=}{{\hbox{<}\hbox{=}}}2 {&&}{{\hbox{\&}\hbox{\&}}}2 {||}{{\hbox{|}\hbox{|}}}2}
\lstdefinestyle{cat}{language=SwiftSheet, basicstyle=\ttfamily\scriptsize,
  aboveskip=1pt, belowskip=2pt}

\makeatletter
\DeclareRobustCommand\op[1]{\texttt{\@tfor\@c:=#1\do{\hbox{\@c}}}}
\makeatother
\newcommand\grp[1]{\par\vspace{1pt}\noindent{\color{sheetBlue}\bfseries\footnotesize #1}\par}
\newcommand\at[1]{\texttt{#1}}
\newcommand\uo{\,{\tiny\itshape\textcolor{sheetRed}{(underscored, unofficial)}}}

\tikzset{
  lib/.style={draw=sheetBlue, fill=sheetBlue!7, font=\ttfamily\scriptsize, minimum height=4.3mm,
              minimum width=46mm, inner sep=1pt, anchor=east},
  lbl/.style={font=\scriptsize, text=black!85, inner sep=1pt, anchor=west, align=left},
  ttl/.style={font=\bfseries\scriptsize},
  call/.style={->, thick, draw=sheetGreen!70!black},
  copy/.style={->, thick, draw=sheetOrange},
  hide/.style={->, thick, dashed, draw=sheetGrey},
}

\begin{document}

\sheettitle{Swift attributes \& compiler directives — a catalogue}{swift · folio}

\oneliner{An \textbf{attribute} (\texttt{@name(args)}) changes how the compiler treats a declaration
or type — what it exports, how it dispatches, what it checks — without changing its code; a
\textbf{directive} (\texttt{\#if}, \texttt{\#warning}, \texttt{\#fileID}…) acts at compile time
on the source itself. The ones that matter most are about \textbf{boundaries}: module/ABI
(\texttt{@inlinable}, \texttt{@frozen}), language (\texttt{@objc}), concurrency
(\texttt{@preconcurrency}) and build configuration (\texttt{\#if}).}

\vspace{2pt}
\noindent\begin{tikzpicture}[sheet]
  \node[ttl, anchor=east] at (5.2,3.2) {library, built with library evolution};
  \node[ttl, anchor=west] at (6.0,3.2) {what a client app compiled against v1 gets — and what v2 may still change};
  \draw[sheetRed, very thick, dashed] (5.6,3.35) -- (5.6,-0.65);
  \node[font=\tiny\bfseries, text=sheetRed, fill=white, inner sep=1pt] at (5.6,3.5) {module boundary = \texttt{.swiftinterface} + ABI};
  \foreach \d/\s/\t [count=\i from 0] in {
     {public func f()}/call/{\textbf{symbol + signature}: the client \emph{calls} it; v2 may rewrite the body freely},
     {@inlinable public func g()}/copy/{body \textbf{copied into the interface}: the client may inline it — v1's body stays in the app until it is rebuilt},
     {@usableFromInline internal func h()}/call/{\textbf{symbol exported, name hidden}: only inlinable code can call it; client source can't},
     {@frozen public struct Point}/copy/{\textbf{layout promised}: stored inline, fields read directly; v2 can never add/reorder stored properties},
     {public struct Config}/hide/{layout \textbf{hidden}: size + fields via runtime metadata and accessors; v2 may add stored properties},
     {@backDeployed(before: iOS 17)}/copy/{a \textbf{fallback copy} is emitted into the client and used when the OS's library is older (5.8)}} {
    \node[lib] (d\i) at (5.2,2.75-\i*0.52) {\d};
    \draw[\s] (d\i.east) -- (6.0,2.75-\i*0.52);
    \node[lbl] at (6.05,2.75-\i*0.52) {\t};
  }
  \node[font=\tiny, anchor=west] at (6.05,-0.45) {\textcolor{sheetGreen!60!black}{\textbf{$\to$} called through a symbol} \quad
     \textcolor{sheetOrange}{\textbf{$\to$} code or layout baked into the client} \quad \textcolor{sheetGrey}{\textbf{- - $\to$} resolved at run time}};
  \node[font=\tiny, anchor=east, align=right] at (5.2,-0.45) {6.3 (SE-0497): \texttt{@export(implementation)} = body for clients,
     no symbol\\promised · \texttt{@export(interface)} = a symbol, body hidden};
\end{tikzpicture}

\begin{multicols}{2}
\raggedright

\grp{Library evolution \& optimisation}
{\par\noindent\scriptsize\setlength{\tabcolsep}{2pt}\renewcommand{\arraystretch}{1.08}%
\begin{tabularx}{\linewidth}{@{}>{\raggedright\arraybackslash}p{26mm}>{\raggedright\arraybackslash}X@{}}
\at{@inlinable} & body is part of the public interface; may use only \texttt{public} or
  \at{@usableFromInline} declarations \\
\at{@usableFromInline} & on \texttt{internal}/\texttt{package}: ABI-visible, source-invisible \\
\at{@frozen} & public struct/enum: layout (or cases) fixed forever; ignored without library evolution \\
\at{@backDeployed(before:)} & ships a copy of a new API into clients on older OSes (Apple SDK-style libraries) \\
\at{@inline(\_\_always)} / \at{(never)} & inlining hint / forbid\uo; \at{@inline(always)} (6.3, SE-0496)
  \emph{errors} if a direct call can't be inlined, implies \at{@inlinable} on public API \\
\end{tabularx}\par}

\grp{API surface}
{\par\noindent\scriptsize\setlength{\tabcolsep}{2pt}\renewcommand{\arraystretch}{1.08}%
\begin{tabularx}{\linewidth}{@{}>{\raggedright\arraybackslash}p{26mm}>{\raggedright\arraybackslash}X@{}}
\at{@discardableResult} & callers may ignore the return value without a warning \\
\at{@available(*, deprecated,} \at{renamed: "load(id:)")} & warning + Xcode fix-it; \at{message:},
  \at{unavailable} (error), \at{introduced:}/\at{obsoleted:}, \at{noasync} \\
\at{@available(iOS 17, *)} & callers must be in an \at{if \#available} block or be gated themselves \\
\at{@main} & the type's \at{static func main()} (may be \at{async throws}) is the entry point;
  \at{@UIApplicationMain} is an \textbf{error} in Swift 6 mode (SE-0383) \\
\at{@testable import M} & tests see \texttt{internal}; needs \at{ENABLE\_TESTABILITY} (Debug) \\
\at{@\_spi(Beta) public} & visible only to \at{@\_spi(Beta) import}\uo \\
\at{package} & access for modules of one package (5.9, SE-0386) \\
\at{: @retroactive P} & you own neither type nor protocol (6.0, SE-0364) \\
\end{tabularx}\par}

\grp{Objective-C interop}
{\par\noindent\scriptsize\setlength{\tabcolsep}{2pt}\renewcommand{\arraystretch}{1.08}%
\begin{tabularx}{\linewidth}{@{}>{\raggedright\arraybackslash}p{26mm}>{\raggedright\arraybackslash}X@{}}
\at{@objc} / \at{@objc(name)} & exposed to the ObjC runtime (selectors, \at{\#selector}, KVC) \\
\at{@objcMembers} & \at{@objc} on every compatible member — incl. extensions and subclasses \\
\at{@nonobjc} & opt one member out \\
\at{@objc dynamic} & message dispatch via \at{objc\_msgSend}: KVO, swizzling; \at{dynamic} alone
  is not KVO-observable \\
\at{@convention(c)} / \at{(block)} & function type with a C / ObjC-block ABI \\
\end{tabularx}\par}

\grp{Concurrency \& closures (details: concurrency sheets)}
{\par\noindent\scriptsize\setlength{\tabcolsep}{2pt}\renewcommand{\arraystretch}{1.08}%
\begin{tabularx}{\linewidth}{@{}>{\raggedright\arraybackslash}p{26mm}>{\raggedright\arraybackslash}X@{}}
\at{@preconcurrency import M} & downgrade Sendable diagnostics from an unmigrated module (5.6) \\
\at{: @preconcurrency P} & conformance whose isolation is checked at \textbf{run time} (6.0, SE-0423) \\
\at{@unchecked Sendable} & ``trust me'': no checking of stored properties — you own the locking \\
\at{@TaskLocal static var id} & value bound per task tree: \at{\$id.withValue("r1") \{ \ldots\ \}};
  child tasks inherit, \at{Task.detached} doesn't \\
\at{@escaping @Sendable} \at{@MainActor () }\op{->}\at{ Void} & closure type: may outlive the call ·
  may run concurrently · runs on the main actor \\
\at{nonisolated} · \at{(unsafe)} & opt a member out of its actor · opt out of checking \\
\at{consuming} / \at{borrowing} & parameter ownership; \at{consume x} ends a lifetime,
  \at{copy x} copies explicitly (5.9) \\
\end{tabularx}\par}

\grp{Type roles, Interface Builder, Core Data}
{\par\noindent\scriptsize\setlength{\tabcolsep}{2pt}\renewcommand{\arraystretch}{1.08}%
\begin{tabularx}{\linewidth}{@{}>{\raggedright\arraybackslash}p{26mm}>{\raggedright\arraybackslash}X@{}}
\at{@propertyWrapper} · \at{@resultBuilder} & the type becomes a wrapper / a DSL builder (own sheets) \\
\at{@dynamicMemberLookup} · \at{@dynamicCallable} & \at{x.name} / \at{x(1, 2)} rewritten to subscripts / calls \\
\at{@IBOutlet} · \at{@IBAction} & storyboard connections; imply \at{@objc} \\
\at{@NSManaged var} & Core Data supplies the accessors at run time \\
\at{@NSCopying var} & the setter stores \at{newValue.copy()} \\
\at{@warn\_unqualified\_access} & warn unless called as \at{self.min()} — avoids the global \at{min} \\
\end{tabularx}\par}

\columnbreak

\grp{Compiler directives}
\begin{lstlisting}[style=cat]
#if os(iOS) || os(visionOS)       // os arch canImport targetEnvironment
#if canImport(UIKit) && !targetEnvironment(macCatalyst)
#if targetEnvironment(simulator)  // simulator | macCatalyst
#if compiler(>=6.0)   // the COMPILER version (any language mode)
#if swift(>=6.0)      // the LANGUAGE MODE (-swift-version 6)
#if hasFeature(ExistentialAny)    // 5.8 SE-0362; hasAttribute(retroactive)
#if DEBUG             // your -D flag (SWIFT_ACTIVE_COMPILATION_CONDITIONS)
Text("Hi")
#if os(iOS)
  .padding()          // 5.5 SE-0308: #if around postfix members
#endif
#warning("TODO: remove")  // #error("...") stops the build
#sourceLocation(file: "gen.swift", line: 1)  // generated code maps back
#fileID    // "App/Feed.swift"; #file = #fileID in Swift 6 mode
#filePath  // full build-machine path - baked into the binary!
#function  // "load(id:)"   #line #column   #dsohandle: image handle
func run(isolation: isolated (any Actor)? = #isolation)  // caller's actor
macro stringify<T>(_ v: T) -> (T, String) =
  #externalMacro(module: "Macros", type: "StringifyMacro")
\end{lstlisting}

\section{Traps}
\begin{itemize}
  \trap{\texttt{swift(}\op{>=}\texttt{6)} is the language \emph{mode}; a Swift 6 compiler in 5
        mode fails it but passes \texttt{compiler(}\op{>=}\texttt{6)}.}
  \trap{Inactive \texttt{\#if} branches must still \emph{parse} — except under
        \texttt{swift()}/\texttt{compiler()} checks.}
  \trap{\texttt{@inlinable} fixes reach clients only when they rebuild; \texttt{@frozen} is forever.}
  \trap{\texttt{@unchecked Sendable} silences the checker; it proves nothing.}
  \trap{\texttt{@Observable}, \texttt{@Model} look like attributes but are \textbf{macros}
        (they expand into code); \texttt{@MainActor} is a global-actor attribute.}
\end{itemize}

\section{Remember}
\textbf{``Attributes move boundaries: module (\texttt{@inlinable}, \texttt{@frozen}), language
(\texttt{@objc}), isolation (\texttt{@preconcurrency}); \texttt{\#if} moves code out of the build.''}


\end{multicols}

\noindent{\footnotesize\color{sheetGrey}\textit{Related:} swift-keypaths-access-control (access
levels) · swift-enums-pattern-matching (\texttt{@frozen} enums) · macros-and-result-builders ·
swift6-strict-concurrency · objc-interop · modularization-spm · swift-syntax-closures-calls}

\end{document}
